\chapter{常见问题}
\markboth{附录 C 常见问题}{附录 C 常见问题}

\begin{description}[style=nextline,leftmargin=0pt]

\item[它和桌面 FreeCAD 是什么关系？]
同一个引擎、同一种文档格式。MobileCAD 重写了界面（Kotlin + Compose），并用 Python 模拟 Qt，让 FreeCAD
自带的任务面板能在触屏上运行。几何算法、文件格式与 FreeCAD 完全一致。

\item[能替代电脑上的 CAD 吗？]
对“设计—修改—出图”的中等复杂度工作，可以。超大装配（上千个零件）、复杂曲面重建、大规模有限元，
桌面版仍更合适。建议的分工：手机/平板做构思、修改、现场测量与审图；桌面做重计算与批量处理。

\item[支持 DWG 吗？]
不支持。DWG 需要外部转换器（ODA/LibreDWG），未打包进 APK。请把 DWG 另存为 \textbf{DXF} 再导入。

\item[有 Python 控制台和插件吗？]
没有交互式控制台，也没有插件管理器。但\textbf{宏仍可运行}（\cmd{Std\_DlgMacroExecute}），
FreeCAD 的 Python 命令也照常工作——BIM/CAM/FEM 的界面本身就是 Python 写的。

\item[文件存在哪里？会自动保存吗？]
位置由你通过系统文件选择器决定，应用不强制目录。自动保存是“崩溃恢复副本”，不是正式保存：
每 60 秒写一份到应用私有目录，异常退出后自动恢复。请仍然养成手动保存的习惯。

\item[能同时打开两个文档吗？]
不能。MobileCAD 是单文档界面，新开文件会替换当前视图。

\item[为什么有些菜单项是灰的？]
那些是桌面窗口管理类命令（多窗口、停靠面板、工具栏定制、宏录制等），移动端没有对应物。点按会提示
“该命令尚未在移动端提供”。

\item[中文界面怎么开启？]
默认跟随系统：系统是中文就显示简体中文。也可在 \menupath{编辑 → 偏好设置 → 语言} 里固定为
\ui{简体中文}、\ui{英语} 或 \ui{系统语言}。

\item[支持哪些设备？]
Android 8.0 及以上、arm64 架构。手机、平板均可；Quest 等头显使用 MobileCAD XR 版本。安装包约 165 MB，
运行时约 1 GB。

\item[能离线使用吗？]
能。建模、装配、出图、有限元、CAM 全部本地完成。只有“检查更新”、打开网页帮助/捐赠链接时才联网。

\item[有限元能算多大？]
取决于设备内存与耐心。手机上建议单元数控制在几十万以内；先用粗网格验证，再局部加密。
只支持 CalculiX 求解器，Netgen 网格器不可用。

\item[能直接打印图纸吗？]
工程图可导出 PDF/SVG/DXF；移动端没有系统打印对话框，“打印”等同于导出当前页 PDF。

\item[GPL 许可证可以商用吗？]
可以。GPL-3.0 允许商业使用，但如果你分发修改后的程序，需要同时提供对应源代码。内部使用不受限制。

\item[怎样反馈问题或贡献翻译？]
在 \url{https://github.com/LucasCoraca/mobilecad} 提 issue（附设备型号、Android 版本、应用版本、
FreeCAD 提交号、复现步骤与截图）。翻译词表在 \file{app/app/src/main/assets/i18n/}，
中文条目可以直接提交 PR。

\end{description}
